\chapter{Assembly of the local certificates}
\label{ch:cover}
\label{component:pe1088004af65}
\label{pe1088004af65:sec:complete-source-routing}

The preceding chapters reduce every valid inequality to events, the
closed faces, and the two contact families C1 and C2. Events and closed
faces already have separated certificates. We now prove the same
assertion for the contact families.

\begin{theorem}[Canonical certificate]
For every source in either family SG10 with an actual finite
noncommuting stationary zero, the complete homogeneous state polynomial
is a finite sum of squares in $K_c\otimes_{\mathbb R}K_t$ at every
independent projective testing pair and every real or complex state.
\end{theorem}

Here and below, a source coefficient times a testing expression denotes
its separated tensor product. With $D_u=1+u^2$, $D_v=1+v^2$, the full
polynomial to be certified is
\[
\begin{split}
 F={}&4\lambda D_u(W+vX)^2+4nD_uD_vX^2
       +4m(vW-X+uvY-uZ)^2\\
 &+4CD_v(uW-Y)^2+4AD_uD_vY^2
       +(4D+2d)(uW+uvX-Y-vZ)^2\\
 &+(4\alpha+2d)D_u(vY-Z)^2
       +(4\beta+2d)D_v(X+uZ)^2-2dD_uD_vZ^2.
\end{split}\tag{LC2}
\]
Its homogeneous form defines the literal projective endpoints.

\begin{proof}
Normalize $d=1$. Every original coefficient can be restored by
multiplying the roots by $\sqrt d\otimes1$. The source parameter
\[
 \delta_{\rm tot}=\frac d{d+\lambda+n+m+C+A+D+\alpha+\beta}
 \tag{LC3}
\]
separates two possibilities. If it is a positive unit, all coefficients
after normalization are bounded. The proof in
Chapter~\ref{ch:cert-2} covers this domain. Otherwise
$\delta_{\rm tot}$ is a positive infinitesimal, and the contact
coordinates from Chapter~\ref{ch:entry} give the source partitions
of Chapters~\ref{ch:c1-short} and \ref{ch:c2-short}.

For reference, the C1 coordinates satisfy
\[
\begin{gathered}
 0<k\le1,\quad \eta,\zeta>0,\quad x=k^2,\\
 r=1-a\eta\in(0,1),\quad\theta=t\zeta\in(0,1),\quad h=1-\theta,\\
 0<a,t<2,\quad E=1-x\eta h>0,\quad Q=1-x\zeta r>0,\\
 M=1+xr-t>0,\quad J=1+xh-a>0,\quad
 e=k\sqrt{\eta/\zeta}>0.
\end{gathered}\tag{LC4}
\]
The C1 partition first separates bounded and infinitesimal $k$, then
the relative source scales $\eta,\zeta,e$. When $\zeta$ is
noninfinitesimal, the product $L=x\zeta$ gives the remaining three
alternatives. Chapter~\ref{ch:c1-short} proves their source inequalities
and constructs the corresponding full quadratic forms. The literal
faces $D=0$ and $\beta=0$ are included wherever the stationary equations
permit them.

For C2, the reflection
$(u,v;W,X,Y,Z)\mapsto(1/v,1/u;Z,X,Y,W)$ exchanges
$(m,D)$ and $(C,\alpha)$. Chapter~\ref{ch:c2-short} uses it to reduce
the source angle types. The case of two weak complementary events
has a separate radial calculation. With one weak event, the directed
source partition distinguishes bounded and unbounded ratios; the
bounded ratios share one full matrix completion. The remaining
unbounded cases and every $n=0$ boundary are proved there. All these
partitions concern coordinates in $K_c$ alone.

Each local identity certifies the full state quadratic form, including
its two first-normal terms and its state Hessian. Positive source
excesses are retained as whole squares. The remaining coefficient
matrix is bounded and has positive-definite real residue, so the
lemma of Chapter~\ref{ch:tensor} gives a finite tensor sum of squares.
The explicit complementary charts cover all independent testing
angles and state ratios.

The sign changes
$(u,Y,Z)\mapsto(-u,-Y,-Z)$ and
$(v,X,Z)\mapsto(-v,-X,-Z)$ preserve LC2. Its homogeneous identities
include zero and infinite half-angles. State normalizations use a
nonzero testing coordinate and are restored by its square. Finally,
real and imaginary state parts add, since the measurements are real.
This proves the theorem.
\end{proof}

To finish the global argument, write an arbitrary valid source as a
nonnegative combination of at most nine extreme sources, and an
arbitrary testing behaviour as a mixture of at most nine real-qubit
behaviours. Events, closed faces and the theorem above certify every
atom pair. Equation M2 combines their squares and preserves the affine
constant. This proves the separated-slack theorem, and hence the finite
semidefinite representation.
